\documentclass[11pt,a4paper]{article}

% --- Core Packages ---
\usepackage[utf8]{inputenc}
\usepackage[T1]{fontenc}
\usepackage{lmodern}
\usepackage[margin=1in]{geometry}
\usepackage{amsmath,amssymb,amsfonts}
\usepackage{graphicx}
\usepackage{booktabs}
\usepackage{microtype}
\usepackage{hyperref}
\usepackage{natbib}
\usepackage{caption}
\usepackage{subcaption}
\usepackage{xcolor}

% --- Hyperref Setup ---
\hypersetup{
    colorlinks=true,
    linkcolor=blue!70!black,
    citecolor=blue!70!black,
    urlcolor=teal!70!black
}

% --- Document Metadata ---
\title{\textbf{Scalable GeoAI for Measuring Micro-Scale Physical Living Standards and Spatial Disparities: Integrating Open Buildings, Standardized Point-of-Interest Data, and Graph Attention Autoencoders}}

\author{
    \textbf{Dimas RD} \\
    Urban Informatics and Spatial Artificial Intelligence Laboratory \\
    \texttt{Target: ISPRS IJGI / PLOS ONE / EarthArXiv}
}

\date{\today}

\begin{document}

\maketitle

\begin{abstract}
Measuring socio-spatial disparities and physical living conditions in rapidly urbanizing regions of the Global South is constrained by spatial data scarcity, administrative boundary aggregation, and the relative normalization of survey-based wealth models. This study presents an unsupervised GeoAI framework that maps physical living standards and identifies micro-scale spatial discontinuities across the 12,800 km$^2$ Jabodetabek metropolitan region and Karawang Regency, Indonesia. We combine Google Open Buildings v3 (2.5D building heights) intersected with European Commission JRC GHS-BUILT-C residential masks, Overture Maps Foundation points of interest (POIs), Sentinel-2 multi-spectral indices, Landsat land surface temperature, Open-Meteo air quality reanalysis, and NASA VIIRS nighttime radiance onto 17,300 Uber H3 Resolution 8 hexagonal cells ($\sim$0.74 km$^2$). We formulate a 10-dimensional Physical Living Standard Index (PLSI) calibrated against international shelter, environmental health, and pedestrian accessibility guidelines (UN-Habitat 9.0 m$^2$ living area baseline, WHO PM$_{2.5}$ guidelines, and 15-minute walking catchments). To detect spatial boundaries without ground-truth training labels, we deploy a Spatial Graph Attention Autoencoder (GAT-AE) over hexagonal geometric adjacency. Constraining latent dimensionality ($d = 2$) encourages local spatial continuity, causing reconstruction error to increase where abrupt feature differences occur across physical and social boundaries dividing master-planned residential developments from dense informal settlements. Attributed error decomposition distinguishes socio-spatial boundaries from natural land-cover transitions. A four-quadrant comparison against Meta's Relative Wealth Index illustrates that asset-based measures can diverge from in-situ physical living conditions in dense, asset-holding informal neighborhoods. This framework provides urban and regional planners with an auditable, continuous instrument to evaluate physical living conditions without survey lag.
\end{abstract}

\noindent\textbf{Keywords:} GeoAI, spatial segregation, informal settlements, Graph Attention Networks, H3 discrete global grid, Overture Maps, urban living standards, Global South, urban planning.

\vspace{1em}
\hrule
\vspace{1.5em}

% --- Sections ---
\section{Introduction}

Rapid urbanization across the Global South has produced complex spatial configurations, where dense informal settlements (such as Latin American \textit{favelas}, South Asian \textit{bustees}, or Southeast Asian \textit{kampungs}) frequently neighbor master-planned residential developments and commercial centers \citep{Dovey2021, Kuffer2016}. In large metropolitan regions such as Greater Jakarta (Jabodetabek), Indonesia---home to over 35 million residents---these adjacent neighborhoods exhibit measurable differences in dwelling volume, open green space, pedestrian connectivity, and environmental exposure. Tracking these micro-scale variations remains challenging for urban and regional planning authorities who need consistent, fine-grained spatial evidence.

\subsection{The Challenge of Spatial Data Scarcity in Developing Regions}
A central constraint in evaluating micro-scale urban disparities is the scarcity of spatial data in many low- and middle-income regions \citep{Abascal2022}. In contrast to high-income countries where fine-grained municipal property records, tax registries, and high-frequency census blocks are maintained, developing metropolises often depend on official censuses conducted at multi-year intervals and aggregated to coarse administrative boundaries (such as districts or regencies). Administrative aggregation introduces the Modifiable Areal Unit Problem (MAUP; \citealt{Openshaw1984}), obscuring local spatial variations that occur across contiguous streets and residential quarters \citep{Reardon2004}.

Crowdsourced spatial data, most notably OpenStreetMap (OSM), have been widely used to supplement public records. However, in many developing cities, volunteered geographic information exhibits spatial unevenness and attribute gaps \citep{BarringtonLeigh2017}. Commercial corridors and central business districts often have relatively dense contributions, whereas peripheral and informal residential areas frequently contain unclassified building footprints (e.g., tagged simply as \texttt{building=yes}) or incomplete point-of-interest (POI) records for basic services like healthcare and primary education.

\subsection{Characteristics of Supervised Machine Learning Wealth Indices}
Recent advances in Earth observation and machine learning have facilitated large-scale estimation of socioeconomic indicators \citep{Burke2021, Chi2022, Jean2016}. A prominent global resource is Meta's Relative Wealth Index (RWI; \citealt{Chi2022}), which predicts household wealth at a 2.4~km grid resolution across 135 low- and middle-income countries using satellite imagery, mobile network connectivity, and topographic data calibrated against Demographic and Health Surveys (DHS).

While valuable for broad regional screening and humanitarian targeting, supervised wealth models have specific characteristics that warrant attention in urban spatial analysis:
\begin{enumerate}
    \item \textbf{Relative Index Normalization}: RWI models are trained on within-country survey asset scores normalized to standard deviations from the national mean (mean zero, unit variance). Consequently, an RWI score reflects relative economic position within a specific country rather than an absolute measure of physical living conditions. A score of $+1.0$ in a predominantly rural country represents a different built environment than $+1.0$ in an urbanized economy, which limits standardized cross-border comparisons.
    \item \textbf{Divergence Between Household Assets and the Physical Built Environment}: Household surveys typically record ownership of consumer durable assets, such as mobile phones, televisions, refrigerators, and personal vehicles \citep{Rutstein2004}. In dense urban informal settlements, households may possess consumer durables while living in environments with limited living floor area ($<9.0\text{ m}^2$ per person), unpaved lanes, restricted drainage, and elevated particulate air pollution. Models trained on durable asset ownership can classify these areas as relatively well-off, reflecting household consumer goods rather than neighborhood physical living conditions.
    \item \textbf{Data Access and Auditability}: Several high-resolution predictive models incorporate proprietary telecommunications or social platform connectivity data \citep{Blumenstock2015}. For public planning agencies, relying on closed data pipelines can complicate independent audits and long-term methodological replication.
\end{enumerate}

\subsection{Research Objectives and Questions}
To complement existing survey-trained wealth indices, this study develops an unsupervised, cloud-native GeoAI framework that maps absolute physical living standards and identifies micro-scale spatial boundaries using open spatial datasets. We structure the inquiry around three research questions:
\begin{itemize}
    \item \textbf{RQ1}: How can multi-source open remote sensing and open spatial geometries be integrated into an unsupervised, continuous Physical Living Standard Index that is comparable across administrative boundaries?
    \item \textbf{RQ2}: Can an unsupervised Graph Attention Autoencoder (GAT-AE) detect micro-scale spatial disparity and segregation boundaries without ground-truth training labels?
    \item \textbf{RQ3}: How do in-situ physical living standards diverge from satellite-predicted consumer asset wealth (Meta RWI) across dense urban residential morphologies?
\end{itemize}

\subsection{Primary Contributions}
This paper provides four main contributions:
\begin{enumerate}
    \item \textbf{Derivation of the Physical Living Standard Index (PLSI)}: We formulate an interpretable index combining 10 physical dimensions (living floor area, built density, pedestrian amenity access, bare soil, nighttime radiance, vegetative canopy, thermal comfort, water drainage, and air quality). The indicators are anchored to recognized physical baselines, including UN-Habitat living space criteria ($9.0\text{ m}^2$ per person), WHO air quality guidelines, and 15-minute walking catchments.
    \item \textbf{Unsupervised Disparity Detection via Graph Attention Autoencoders}: We deploy a Spatial GAT Autoencoder over a discrete global grid (Uber H3 Resolution 8, $\sim$0.74~km$^2$). By constraining latent dimensionality ($d = 2$) across geometric adjacency, the model establishes local continuity expectations, causing reconstruction error to increase where abrupt feature transitions occur. Attributed error decomposition separates socio-spatial discontinuities from biophysical land-cover edges.
    \item \textbf{Open Data Integration Pipeline}: We demonstrate a reproducible pipeline combining Google Open Buildings v3 (2.5D heights), Overture Maps Foundation standardized point-of-interest data, and open satellite observations (Sentinel-2, Landsat, VIIRS), reducing dependence on proprietary datasets.
    \item \textbf{Planning Relevance and Divergence Analysis}: We apply the Spatial Palma Ratio and a four-quadrant comparison framework (Asset Wealth vs. Physical Infrastructure) to provide urban and regional planners with an auditable tool for identifying infrastructure deficits in dense informal areas.
\end{enumerate}

\subsection{Structure of the Paper}
The remainder of this article is organized as follows: Section 2 reviews relevant literature on remote sensing wealth estimation, 3D building morphology, pedestrian accessibility, and graph neural networks, concluding with the four-quadrant divergence framework. Section 3 outlines the study area and cloud data ingestion workflow. Section 4 presents the 10-dimensional feature engineering methodology and Spatial GAT-AE architecture. Section 5 discusses empirical results and validation benchmarks. Section 6 addresses planning implications and data limitations. Section 7 concludes with considerations for global transferability.

\section{Literature Review and Theoretical Framework}

Evaluating urban living conditions across expanding metropolitan regions in the developing world draws on four complementary fields: remote sensing of socioeconomic indicators, 3D building morphology, pedestrian accessibility modeling, and graph neural networks for spatial anomaly detection. We synthesize these areas into a four-quadrant divergence framework that contrasts asset-based wealth with in-situ physical living standards.

\subsection{Remote Sensing and Machine Learning for Socioeconomic Estimation}
Satellite Earth observation has long supported socioeconomic monitoring. Nocturnal luminosity (DMSP and VIIRS) serves as an established proxy for regional economic activity and electrification \citep{Henderson2012, Mellander2015}, but encounters saturation in commercial cores and blooming across unpopulated surfaces \citep{Zhao2019}. To improve spatial resolution, transfer learning models use daytime multi-spectral imagery to predict nighttime radiance, learning visual features that correlate with local survey expenditures \citep{Burke2021, Jean2016, Yeh2020}. Meta's Relative Wealth Index (RWI; \citealt{Chi2022}) scaled this approach to 135 countries at 2.4~km resolution by training gradient-boosted trees on satellite, mobility, and connectivity data against Demographic and Health Surveys (DHS). 

However, survey asset targets capture household durable goods ownership \citep{Rutstein2004}. In dense urban settlements, asset accumulation can decouple from physical living conditions: households may own mobile devices and appliances while residing in overcrowded structures with unpaved access and open drainage canals.

\subsection{Built-Environment Morphology and Living Space Sufficiency}
Urban remote sensing has traditionally identified informal settlements through distinct morphological signatures: dense building configurations, organic street layouts, small footprints, and heterogeneous roofing \citep{Kohli2012, Sliuzas2010}. Early studies framed detection as binary ``slum vs. formal'' image classification \citep{Kuffer2016, Kuffer2020, Wurm2019}. While computationally effective, binary classifications obscure progressive consolidation \citep{Dovey2020} and ignore the multidimensional continuum of living conditions \citep{Abascal2022, Lilford2019}.

To measure conditions continuously, UN-Habitat \citep{UNHabitat2003, UNHabitat2020} establishes five shelter deprivation criteria: durable construction, access to improved water, adequate sanitation, tenure security, and sufficient living area ($<9.0\text{ m}^2$ per person). Estimating floor area per capita historically proved difficult from 2D overhead imagery. Recent open planetary-scale building datasets, particularly Google Open Buildings v3 \citep{Sirko2021} with 2.5D heights, now enable volumetric morphology estimation across the Global South \citep{Biljecki2023, Biljecki2024}. Intersecting building volumes with ambient population data (WorldPop; \citealt{Tatem2017}) and functional residential masks (JRC GHS-BUILT-C; \citealt{Schiavina2022}) enables continuous estimation of living floor area per capita across regional extents.

\subsection{Pedestrian Accessibility and Open Spatial Data}
Physical living standards also depend on spatial access to basic public amenities, including healthcare, education, retail, and green space \citep{Geurs2004, Handy1997}. The ``15-minute city'' concept emphasizes reachable daily services within a 15-minute walk \citep{Moreno2021}. Cumulative opportunity thresholds (destinations within a fixed walking catchment at 4.0~km/h) offer intuitive, policy-relevant metrics compared to gravity models \citep{Hosford2022, Weiss2018}.

In fragmented cities, Euclidean buffer distances can substantially overestimate accessibility \citep{Boeing2017}. Rapid suburban expansion frequently situates enclosed master-planned developments beside informal settlements \citep{Dovey2021, Firman2004}. Perimeter walls and highway barriers require multi-kilometer pedestrian detours, making topological network routing essential. While OpenStreetMap enables open routing \citep{Boeing2017}, volunteer contributions remain uneven in peripheral areas \citep{BarringtonLeigh2017}. The Overture Maps Foundation \citep{Overture2023, Overture2024} addresses this gap by combining open registries and algorithmic extraction into a standardized Global Entity Reference System (GERS).

\subsection{Graph Neural Networks and Spatial Anomaly Detection}
Unsupervised disparity detection across continuous space requires structural learning methods beyond linear spatial autocorrelation statistics \citep{Anselin1995, Reardon2004}. Graph Neural Networks (GNNs), such as Graph Convolutional Networks (GCNs; \citealt{Kipf2016}) and Graph Attention Networks (GATs; \citealt{Velickovic2018}), learn representations over non-Euclidean spatial graphs where hexagonal cells form nodes and geometric adjacency forms edges \citep{Zhang2022}.

Deep GNNs often encounter ``over-smoothing,'' where repeated neighborhood message passing homogenizes node embeddings \citep{Chen2020, Kipf2016}. In unsupervised anomaly detection, this characteristic can be repurposed. By training a Graph Attention Autoencoder (GAT-AE) with a narrow bottleneck ($d=2$) over pure geometric adjacency, the model enforces local spatial continuity. Where an edge crosses an abrupt socio-spatial boundary, message passing aggregates conflicting feature distributions, resulting in elevated reconstruction errors at the interface. Drawing on attributed graph anomaly frameworks (DOMINANT; \citealt{Ding2019}), attributed error decomposition separates built-environment indicators from biophysical land-cover edges:
\begin{equation}
\text{Reconstruction Error}_i = \sum_{k=1}^M (x_{ik} - \hat{x}_{ik})^2
\end{equation}

\subsection{Theoretical Synthesis: The 4-Quadrant Divergence Framework}
We synthesize these domains into the \textbf{4-Quadrant Divergence Framework}, conceptualizing urban space along two axes: Consumer Asset Wealth (Meta RWI) and In-Situ Physical Living Standards (PLSI).

\begin{figure}[htbp]
\centering
\begin{verbatim}
                  High Physical Living Standard (PLSI)
                                  ▲
                                  │
                 Quadrant III     │     Quadrant I
            (Subsidized Formal /  │ (Integrated Affluent /
             Well-Serviced Low-   │  Master-Planned Enclaves)
             Income Peripheries)  │
                                  │
◄─────────────────────────────────┼─────────────────────────────────►
Low Asset Wealth (Meta RWI)       │       High Asset Wealth (Meta RWI)
                                  │
                 Quadrant II      │     Quadrant IV
            (Unserviced Informal  │  (Asset-Dense, Infrastructure-
             Deprivation / Rural  │   Deprived Informal Settlements)
             Poverty)             │
                                  ▼
                   Low Physical Living Standard (PLSI)
\end{verbatim}
\caption{The 4-Quadrant Divergence Framework contrasting asset-based wealth with in-situ physical living standards.}
\label{fig:four_quadrant_framework}
\end{figure}

\begin{itemize}
    \item \textbf{Quadrant I (High Asset, High Physical)}: Formal residential communities with large dwelling areas, tree cover, and walkable amenities.
    \item \textbf{Quadrant II (Low Asset, Low Physical)}: Unplanned settlements and impoverished peri-urban areas lacking private assets and public services.
    \item \textbf{Quadrant III (Low Asset, High Physical)}: Subsidized public housing with modest household assets but planned dwelling volumes, drainage, and paved access.
    \item \textbf{Quadrant IV (High Asset, Low Physical)}: Established dense informal settlements where households hold durable goods (appliances, vehicles) but endure floor-space deprivation ($<9.0\text{ m}^2$), open drainage, and elevated particulate air pollution.
\end{itemize}

Supervised models trained on survey asset indices often conflate Quadrant IV with Quadrant I, obscuring local infrastructure deficits. The unsupervised PLSI framework provides an empirical tool to evaluate both dimensions independently.


% --- Bibliography ---
\newpage
\bibliographystyle{apalike}
\bibliography{references}

\end{document}
